\begin{explanationbox}
	\textbf{\textcolor{CExplanation}{一句话直觉}}
	正弦型函数的图像由正弦曲线经伸缩、平移得到，周期由 $\omega$ 唯一确定，单调区间只需将 $\omega x+\varphi$ 整体代入 $\sin$ 的单调区间，再解出 $x$。

	\medskip
	\textbf{\textcolor{CExplanation}{核心拆解}}
	\begin{description}[style=nextline, leftmargin=2.8em, labelsep=0.5em, itemsep=0.45em]
		\item[\textbf{要点一（周期与 $\omega$）：}]
			周期 $T = \frac{2\pi}{\omega}$，$\omega$ 越大，周期越短.
		\item[\textbf{要点二（整体代换法）：}]
			设 $u = \omega x + \varphi$，则单调性由 $u$ 所在区间决定；列出 $u$ 满足的不等式，代回 $x$.
	\end{description}

	\medskip
	\textbf{\textcolor{CExplanation}{几何本质（图像伸缩平移）}}
	$y=A\sin x$ 的图像横坐标变为原来的 $\frac{1}{\omega}$ 倍得到 $y=A\sin(\omega x)$ 的简图；再向左平移 $\frac{\varphi}{\omega}$（$\varphi > 0$ 向左）得最终图像。周期变为原来的 $\frac{1}{\omega}$ 倍，单调区间随之缩短或伸长。

	\medskip
	\textbf{\textcolor{CExplanation}{代数意义（整体代换解不等式）}}
	原问题转化为：已知 $\sin u$ 的单调区间，求 $u = \omega x + \varphi$ 中 $x$ 的范围，即解一元一次不等式组。

	\medskip
	\textbf{\textcolor{CExplanation}{考点价值（高频题型）}}
	\begin{itemize}[leftmargin=*, itemsep=0.5em, topsep=0.4em]
		\item \textbf{考法一（求周期）：} 直接代入公式 $T = \frac{2\pi}{\omega}$（针对 $\omega>0$）.
		\item \textbf{考法二（求单调区间）：} 列出区间不等式，解出 $x$（注意整数 $k$）.
		\item \textbf{考法三（求参数范围）：} 已知单调区间反向确定 $\omega$ 或 $\varphi$.
	\end{itemize}

	\medskip
	\textbf{\textcolor{CExplanation}{顿悟点}}
	\textbf{关键在“整体代换”——把 $\omega x+\varphi$ 当作一个角，而不是分开看 $x$ 和相位。}

	\medskip
	\textbf{\textcolor{CExplanation}{使用场景}}
	\begin{itemize}[leftmargin=*, itemsep=0.5em, topsep=0.4em]
		\item \textbf{场景一（纯正弦型函数）：} 直接给出 $y = A\sin(\omega x + \varphi)$，求周期与单调区间.
		\item \textbf{场景二（含参数问题）：} 已知函数的单调区间特征，求 $\omega$ 或 $\varphi$ 的取值范围.
	\end{itemize}

\end{explanationbox}
